% Conceptual relationships, not a temporal pipeline or measured gains.
% Final width is below 169 mm. All outgoing requests share the east-side port.
\begin{tikzpicture}[
  x=1mm,y=1mm,
  font=\figurefont\small,
  every node/.style={text=BookInk},
  source/.style={
    rectangle,draw=BookTeal,fill=BookTeal!8,
    line width=1pt,text width=39mm,minimum height=22mm,
    inner sep=2mm,align=center
  },
  request/.style={
    rectangle,draw=BookBlue,fill=BookPaper,
    line width=0.8pt,text width=55mm,minimum height=12mm,
    inner sep=1.5mm,align=center
  },
  link/.style={draw=BookTeal,line width=0.9pt},
  arrow/.style={link,-{Latex[length=2mm,width=1.4mm]}},
  change/.style={font=\figurefont\footnotesize,anchor=south,text=BookMuted}
]
  \node[source] (system) at (23,0) {持续运行的\\花卉识别系统};
  \node[font=\figurefont\footnotesize,align=center,text width=42mm]
    at (23,-18) {已有表示、模型与记忆\\附带机构问答服务};

  \node[request] (environment) at (124,27) {环境随季节/地区漂移\\光照、背景与类别比例变化};
  \node[request] (newclass) at (124,9) {类别陆续增加\\学新物种且不遗忘旧物种};
  \node[request] (sensor) at (124,-9) {观测设备更换\\RGB 换成多光谱或红外};
  \node[request] (editing) at (124,-27) {事实修正\\负责人交接与有效日期};

  \coordinate (fork) at (61,0);
  \draw[link] (system.east) -- (fork);
  \draw[link] (61,-27) -- (61,27);
  \foreach \target/\height/\object in {
    environment/27/观测分布,
    newclass/9/顺序学习与旧类保持,
    sensor/-9/观测含义与接入,
    editing/-27/事实有效范围
  } {
    \draw[arrow] (61,\height) -- (\target.west);
    \node[change] at (80,\height+1) {\object};
  }
  \node[font=\figurefont\footnotesize,text=BookMuted,anchor=north]
    at (95,-37) {各类变化可以同时发生};
\end{tikzpicture}
